% Part: methods
% Chapter: induction
% Section: inductive-definitions

\documentclass[../../../include/open-logic-section]{subfiles}

\begin{document}

\olfileid{mth}{ind}{idf}

\olsection{归纳定义}

在逻辑中，我们经常\emph{归纳地}定义各类对象，亦即，通过规定何种对象算作该类对象的规则来定义，这些规则说明了如何从已有的该类对象生成新的该类对象。例如，我们经常通过归纳来定义特定的符号序列，比如一种语言的项和公式。来看一个简单的例子：考虑由字母 $\mathrm{a}$、$\mathrm{b}$、$\mathrm{c}$、$\mathrm{d}$、符号~$\circ$ 以及方括号 $[$ 和 $]$ 组成的符号串，例如 ``$[[\mathrm{c} \circ \mathrm{d}][$''、``$[\mathrm{a}[]\circ]$''、``$\mathrm{a}$'' 或 ``$[[\mathrm{a} \circ \mathrm{b}]\circ
\mathrm{d}]$''。你可能会觉得前两个符号串有些「不对劲」：第一个的括号完全「不平衡」，而第二个可能会让你觉得 ``$\circ$'' 应该「连接」自身有意义的表达式。第三和第四个符号串看起来就好多了：只要出现了 ``$[$''，就一定有一个闭合的 ``$]$'' 与之对应，并且对于任何一个 $\circ$，我们都能在它两边找到被一对括号括起来的「好」表达式。

我们想要精确地规定什么算作一个「好项」。首先，每个字母本身都是好项。任何不仅仅由单个字母构成的对象，都应当具有 ``$[t \circ s]$'' 的形式，其中 $s$ 和 $t$ 本身也是好项。反过来，如果 $t$ 和 $s$ 是好项，那么我们可以在它们中间放一个 $\circ$ 并用一对括号将它们括起来，从而形成一个新的好项。我们可以利用这些操作来\emph{定义}好项的集合。这就是一个\emph{归纳定义}。

\begin{defn}[好项]
  \emph{好项}的集合归纳定义如下：
  \begin{enumerate}
  \item 任何字母 $\mathrm{a}$、$\mathrm{b}$、$\mathrm{c}$、
    $\mathrm{d}$ 都是好项。
  \item 若 $s_1$ 和 $s_2$ 是好项，那么 $[s_1 \circ s_2]$ 也是好项。
  \item 没有其他东西是好项。
  \end{enumerate}
\end{defn}

此定义告诉我们，一个对象算作好项，当且仅当它能根据条款（1）和~(2)
在有限步内构造出来。第一步，我们构造出所有仅由字母本身组成的好项，即
\[
\mathrm{a}, \mathrm{b}, \mathrm{c}, \mathrm{d}
\]
第二步，我们将（2）应用到已经构造出的项上。就会得到
\[
[\mathrm{a} \circ \mathrm{a}], [\mathrm{a} \circ \mathrm{b}],
[\mathrm{b} \circ \mathrm{a}], \dots, [\mathrm{d} \circ \mathrm{d}]
\]
这涵盖了所有两个字母的组合。第三步，我们再次将（2）应用到目前构造出的任意两个好项上。我们会得到新的好项，例如 $[\mathrm{a} \circ [\mathrm{a} \circ
    \mathrm{a}]]$——其中 $t$ 是来自第一步的 $\mathrm{a}$，而 $s$ 是
来自第二步的 $[\mathrm{a} \circ \mathrm{a}]$——以及由第二步中的两个项 $[\mathrm{b} \circ \mathrm{c}]$ 和 $[\mathrm{d}
  \circ \mathrm{b}]$ 构造出的 $[[\mathrm{b} \circ
    \mathrm{c}] \circ [\mathrm{d} \circ \mathrm{b}]]$。如此等等。条款（3）排除了任何未以这种方式构造出来的对象被算作好项的可能。

注意，我们尚未证明每一个感觉上是好的符号串根据这个定义都是好项。然而，应该很清楚，我们能构造出的所有东西实际上感觉上都是好的：括号是平衡的，而且 $\circ$ 连接的部分本身也是好的。

归纳定义的关键特征是，如果你想证明关于所有好项的某种性质，该定义就会指明你必须考虑哪些情况。例如，如果已知 $t$ 是一个好项，那么归纳定义就告诉你 $t$ 可能的样子：$t$ 可以是一个字母，或者对于某一对好项 $s_1$ 和~$s_2$，它可以是 $[s_1 \circ s_2]$。由于条款（3），这些是仅有的可能性。

在证明关于某个归纳定义集合全体元素的断言时，强归纳法变得尤为重要。例如，假设我们要证明对于每个长度为~$n$ 的好项，其中 $[$ 的个数 $< n/2$。这可以被看作一个关于所有~$n$ 的断言：对于每个 $n$，任何长度为~$n$ 的好项中 $[$ 的个数都 $< n/2$。

\begin{prop}
  对任意 $n$，长度为~$n$ 的好项中 $[$ 的个数都 $< n/2$。
\end{prop}

\begin{proof}
为了用（强）归纳法证明这个结果，我们需要证明以下条件命题为真：
\begin{quote}
  如果对每个 $l < k$，任意长度为~$l$ 的好项中 $[$ 的个数都 $< l/2$，
  那么任意长度为~$k$ 的好项中 $[$ 的个数都 $< k/2$。
\end{quote}
为证明这个条件命题，假设其前件为真，即假设对任意 $l<k$，长度为~$l$ 的好项中 $[$ 的个数 $< l/2$。
我们称此假设为归纳假设。我们要证明该结论对长度为~$k$ 的好项同样成立。

于是设 $t$ 是一个长度为~$k$ 的好项。由于好项是归纳定义的，我们有两种情况：
(1)~$t$ 本身是一个字母，或者 (2)~$t$ 是 $[s_1 \circ s_2]$，其中 $s_1$ 和~$s_2$ 是好项。
\begin{enumerate}
\item $t$ 是一个字母。那么 $k = 1$，并且 $t$ 中 $[$ 的个数为~$0$。由于 $0 < 1/2$，结论成立。
\item $t$ 是 $[s_1 \circ s_2]$，其中 $s_1$ 和~$s_2$ 是好项。
  令 $l_1$ 为~$s_1$ 的长度，$l_2$ 为~$s_2$ 的长度。
  那么 $t$ 的长度~$k$ 为 $l_1+l_2+3$（即 $s_1$ 和 $s_2$ 的长度加上三个符号 $[$、$\circ$、$]$）。
  由于 $l_1+l_2+3$ 总是大于 $l_1$，所以 $l_1 < k$。类似地，$l_2 < k$。
  这意味着归纳假设适用于项 $s_1$ 和~$s_2$：$s_1$ 中 $[$ 的个数~$m_1$ 满足 $m_1 < l_1/2$，
  并且~$s_2$ 中 $[$ 的个数~$m_2$ 满足 $m_2 < l_2/2$。

  $t$ 中 $[$ 的个数等于~$s_1$ 中 $[$ 的个数加上~$s_2$ 中 $[$ 的个数再加~$1$，
  即 $m_1 + m_2 + 1$。由于 $m_1 < l_1/2$ 且 $m_2 < l_2/2$，我们有：
  \[
  m_1 + m_2 + 1 < \frac{l_1}{2} + \frac{l_2}{2} + 1 = \frac{l_1+l_2+2}{2} < \frac{l_1+l_2+3}{2} = k/2.
  \]
\end{enumerate}
在每种情况下，我们都证明了（在归纳假设的基础上）$t$ 中 $[$ 的个数 $< k/2$。
由强归纳法，命题得证。
\end{proof}

\begin{prob}
  通过以下方式定义超级好项的集合：
  \begin{enumerate}
  \item 任意字母 $\mathrm{a}$、$\mathrm{b}$、$\mathrm{c}$、$\mathrm{d}$ 都是超级好项。
  \item 如果 $s$ 是超级好项，那么 $[s]$ 也是超级好项。
  \item 如果 $s_1$ 和 $s_2$ 是超级好项，那么 $[s_1 \circ s_2]$ 也是超级好项。
  \item 没有其他东西是超级好项。
  \end{enumerate}
  证明：对长度为~$n$ 的任意超级好项~$t$，其中 $[$ 的个数 $\le n/2 +1$。
\end{prob}

\end{document}